% Exercise Sheet 01
% Define each exercise as a command so it can be placed in any chapter.
%
% Suggested naming scheme:
%   \sheetoneexone
%   \sheetoneextwo
%   \sheetoneexthree
%
% Then call the relevant commands inside:
%   content/exercises/01_symmetries.tex
%   content/exercises/02_basic_group_theory.tex

\NewDocumentCommand{\sheetoneexone}{}{
    \ecs{The symmetric group \(S_n\)}{
        We have defined the symmetric group \(S_n\) as the group of permutations
        on the set \(\{1,2,\dots,n\}\). A permutation \(\sigma\) acts on elements
        \(i \in \{1,2,\dots,n\}\) by
        \[
            i \mapsto \sigma(i).
        \]

        We also denote this permutation as
        \[
            \sigma =
            \begin{pmatrix}
                1 & 2 & \dots & n \\
                \sigma(1) & \sigma(2) & \dots & \sigma(n)
            \end{pmatrix}.
        \]

        For \(1 \leq k \leq n\), let \(a_1,\dots,a_k \in \{1,2,\dots,n\}\) be
        pairwise different. A cycle \((a_1,\dots,a_k)\) is the cyclic permutation
        \(a_1 \mapsto a_2\), \(a_2 \mapsto a_3\), \dots, \(a_k \mapsto a_1\);
        it acts as the identity on all other elements. We call \(k\) the length
        of the cycle.
    }{
        \ecpart{Show that any permutation \(\sigma\) can be written as a product of
        disjoint cycles.}{
            We show this by explicit construction. For any element \(i\in\{1,2,\dots,n\}\) we can construct a cycle by repeatedly applying \(\sigma\). Since the set \(\{1,2,\dots,n\}\) is finite, there exists a smallest \(k\) s.t. \(\sigma^k(i) = i\). We thus constructed a cycle \((i, \sigma(i), \sigma^2(i), \dots, \sigma^{k-1}(i))\). We can repeat this procedure for any element that has not yet been included in a cycle until we have exhausted all elements. By construction, the resulting cycles are disjoint and their product is \(\sigma\).
        }

        \ecpart{A transposition is a cycle of length \(2\). For \(n \geq 2\), show
        that every cycle, and thus every permutation, can be written as a product
        of transpositions. This means that \(S_n\) is generated by transpositions.}{
            This is shown by induction on the length \(k\) of the cycle.
            \begin{itemize}
                \item For \(k=2\) the cycle is already a transposition.
                \item For \(k = 3\), we can write \((a_1,a_2,a_3) = (a_1,a_3)\circ(a_1,a_2)\).
                \item For \(k > 3\), we can write \((a_1,a_2,\dots,a_k) = (a_1,a_k)\circ(a_1,a_2,\dots,a_{k-1})\). By the induction hypothesis, the cycle \((a_1,a_2,\dots,a_{k-1})\) can be written as a product of transpositions, and thus so can \((a_1,a_2,\dots,a_k)\).
            \end{itemize}
        }

        \ecpart{Give an example that shows that the decomposition of a cycle into
        transpositions is not necessarily unique.}{
            For example, the cycle \((1,2,3)\) can be written as \((1,3)\circ(1,2)\) or as \((2,3)\circ(1,3)\).
        }

        \ecpart{Denote by \(\tau_i\) the cycle \((i,i+1)\) for \(1 \leq i < n\).
        Show that they satisfy the relations
        \[
            \tau_i^2 = 1,\quad
            \tau_{i+1}\tau_i\tau_{i+1} = \tau_i\tau_{i+1}\tau_i,\quad \text{and}\quad
            \tau_i\tau_j = \tau_j\tau_i \text{ for } 2 \leq |i-j|.
        \]}{
            \begin{itemize}
                \item \( (i,i+1)\circ(i,i+1)= id \)
                \item \(\tau_{i+1}\tau_i\tau_{i+1}\)
            \end{itemize}
        }

        \ecpart{If \(\sigma\) is decomposed into disjoint cycles of length \(k_i\)
        for \(i \leq n\), define the sign of \(\sigma\) as
        \[
            \operatorname{sign}(\sigma) = (-1)^{N(\sigma)},
            \qquad
            \text{where }
            N(\sigma) = \sum_i (k_i - 1).
        \]
        Using the fact that two products of \(\tau_i\) represent the same
        permutation if and only if they are related by the relations from d),
        show that sign is a group homomorphism from \(S_n\) to \((\{\pm 1\},\cdot)\).}{
            Add your answer here.
        }
    }
}

\NewDocumentCommand{\sheetoneextwo}{}{
    \ec{Chinese remainder theorem}{
        Show that if \(a\) and \(b\) are coprime positive integers, then
        \[
            \mathbb{Z}/a\mathbb{Z} \times \mathbb{Z}/b\mathbb{Z}
        \]
        is isomorphic to
        \[
            \mathbb{Z}/ab\mathbb{Z}.
        \]
        How can you use this result to decompose \(\mathbb{Z}/n\mathbb{Z}\) for
        an arbitrary \(n \in \mathbb{N}\)?
    }{

    }
}

\NewDocumentCommand{\sheetoneexthree}{}{
    \ec{Automorphisms of \(\mathbb{Z}_n\)}{
        Let \(G\) be a group. Show that the set \(\operatorname{Aut}(G)\) of
        invertible group homomorphisms \(G \to G\) is a group. For \(n \in \mathbb{N}\)
        compute the group \(\operatorname{Aut}(\mathbb{Z}_n)\).
    }{
        \begin{itemize}
            \item The identity map is an invertible group homomorphism, so \(id\in\operatorname{Aut}(G)\). Furthermore \(g^{-1}\circ g = id\)
            \item The composition of two invertible group homomorphisms is again an invertible group homomorphism, so \(\operatorname{Aut}(G)\) is closed under composition.
            \item The composition of group homomorphisms is associative.
        \end{itemize}
        There for \(\left(\operatorname{Aut}(G),\circ \right)\) is a group.
        \\
        Let \(f:\mathbb{Z}_n \to \mathbb{Z}_n\) be a homomorphism. Then \( f([i+j]) = [f(i)+f(j)] \) and thus for any \([i]\in\mathbb{Z}_n\) we have \(f([i])=f(i\cdot[1])=i\cdot f([1])\). Thus \(f\) is determined by \(f([1])=[k]\) with \(f([i])=[k\cdot i]\)
        for some \(k\in\{0,1,\dots,n\}\). 
        \(f\) is invertible iff \([k\cdot i] = [k\cdot j]\) implies \([i] = [j]\). Suppose \( \exists d : k=d\cdot\tilde{k} \text{ and } n = d\cdot\tilde{n}\). Then \(f([\tilde{n}]) = [d\cdot\tilde{k}\cdot\tilde{n}] = \tilde{k} \cdot [n]\). Thus \(f([\tilde{n}]) = f([n])\) but \([\tilde{n}] \neq [n] \). Thus for \(f\) to be invertible, \(k\) and \(n\) must be coprime.
        }
}

\NewDocumentCommand{\sheetoneexfour}{}{
    \ecs{The Lorentz group}{
        Recall that the Lorentz group \(O(3,1)\) is the group of linear
        transformations
        \[
            \Lambda : \mathbb{R}^4 \to \mathbb{R}^4
        \]
        that preserve the inner product
        \[
            -x_0^2 + x_1^2 + x_2^2 + x_3^2.
        \]
        We say that \(\Lambda \in O(3,1)\) is \textit{orthochronous} if the first
        component of \(\Lambda(e_0)\) is positive. Otherwise, we say that it is
        \textit{time reversing}. \(\Lambda \in O(3,1)\) is \textit{parity reversing}
        if it reverses orientation and is orthochronous.
    }{
        \ecpart{Show that there is a group homomorphism
        \[
            F : O(3,1) \to \mathbb{Z}_2 \times \mathbb{Z}_2
        \]
        such that the time and parity reversing elements are sent to
        \((1,0)\) and \((0,1)\), respectively.}{
            Add your answer here.
        }

        \ecpart{Show that \(F\) has a section, i.e. that there exists a group
        homomorphism
        \[
            s : \mathbb{Z}_2 \times \mathbb{Z}_2 \to O(3,1)
        \]
        such that \(F \circ s = \mathrm{id}\).}{
            Add your answer here.
        }

        \ecpart{The kernel of \(F\) is called the proper orthochronous Lorentz
        group \(SO(3,1)^+\). Work out, recall, or look up a description of the
        elements of \(SO(3,1)^+\) in terms of boosts and rotations. Show that
        \[
            O(3,1) \cong SO(3,1)^+ \times (\mathbb{Z}_2 \times \mathbb{Z}_2)
        \]
        as sets, not necessarily as groups. Is this also an isomorphism of groups?}{
            Add your answer here.
        }

        \ecpart{Show that conjugation by \(s(-)\) induces a group homomorphism
        \[
            \varphi : \mathbb{Z}_2 \times \mathbb{Z}_2 \to \operatorname{Aut}(SO(3,1)^+)
        \]
        from \(\mathbb{Z}_2 \times \mathbb{Z}_2\) to the group
        \(\operatorname{Aut}(SO(3,1)^+)\) of automorphisms of \(SO(3,1)^+\).}{
            Add your answer here.
        }

        \ecpart{Equip the set \(SO(3,1)^+ \times (\mathbb{Z}_2 \times \mathbb{Z}_2)\)
        with the group structure
        \[
            (\Lambda,(a,b)) \cdot (\Lambda',(a',b'))
            =
            (\Lambda\, \varphi(a,b)(\Lambda'), (a+a', b+b')).
        \]
        Show that this group is isomorphic to \(O(3,1)\).}{
            Add your answer here.
        }
    }
}
